% Conclusion
% Briefly summarize: problem, approach, findings, contributions
% End with a forward-looking statement

\label{sec:conclusion}

This study investigated how Wikipedia’s collaborative governance
shapes neutrality in Talk page discussions. Using sentiment and
toxicity analysis across multiple topics, we found that most comments
were classified as Neutral and showed low toxicity, suggesting that
Wikipedia’s norms help support constructive dialogue even around
controversial subjects.

At the same time, many Talk pages were driven by only a few active
editors, indicating that participation is not always broadly
distributed. This raises questions about how democratic Wikipedia’s
deliberation processes truly are, despite the platform’s open nature.

Our work provides a simple framework for evaluating discourse tone
and participation patterns using machine learning tools. Future work
should examine editor networks, expand to other languages, and
connect Talk page discussions to the actual edits made to articles.
Understanding these dynamics is important for assessing how online
collaborative systems balance openness, neutrality, and shared
decision-making.